\documentclass[11pt,letterpaper]{article}

% ── Geometry ──────────────────────────────────────────────────────────────────
\usepackage[margin=1in]{geometry}

% ── Pagination ────────────────────────────────────────────────────────────────
\usepackage[all]{nowidow}
\brokenpenalty=10000
\raggedbottom

% ── Fonts ─────────────────────────────────────────────────────────────────────
\usepackage[T1]{fontenc}
\usepackage{newpxtext}     % Palatino-based serif (TeX Gyre Pagella)
\usepackage{newpxmath}     % Matching math font

% ── Colors ────────────────────────────────────────────────────────────────────
\usepackage[dvipsnames]{xcolor}
\definecolor{SectionBlue}{HTML}{1B3A5C}
\definecolor{AccentGray}{HTML}{4A4A4A}
\definecolor{QuoteBorder}{HTML}{8B9DAF}
\definecolor{BoxBg}{HTML}{F7F8FA}

% ── Boxes ─────────────────────────────────────────────────────────────────────
\usepackage[framemethod=tikz]{mdframed}

% ── Lists ─────────────────────────────────────────────────────────────────────
\usepackage{enumitem}

% ── Hyperlinks ────────────────────────────────────────────────────────────────
\usepackage{hyperref}
\hypersetup{
  colorlinks=true,
  linkcolor=SectionBlue,
  urlcolor=SectionBlue,
  pdfauthor={punt-labs},
  pdftitle={biff --- Team Communication for Engineers Who Never Leave the Terminal},
  pdfsubject={Product Launch},
}

% ── Section styling ──────────────────────────────────────────────────────────
\usepackage{titlesec}
\titleformat{\section}
  {\Large\bfseries\color{SectionBlue}}
  {}
  {0pt}
  {}

% ── Header/footer ────────────────────────────────────────────────────────────
\usepackage{fancyhdr}
\pagestyle{fancy}
\fancyhf{}
\fancyhead[R]{\footnotesize\textcolor{AccentGray}{\thepage}}
\renewcommand{\headrulewidth}{0pt}

% ══════════════════════════════════════════════════════════════════════════════
% Custom environments
% ══════════════════════════════════════════════════════════════════════════════

% ── \prsection{title} — styled press release section header ──────────────────
\newcommand{\prsection}[1]{%
  \vspace{1em}%
  {\large\bfseries\color{SectionBlue}#1}%
  \vspace{0.3em}\par%
}

% ── customerquote — indented, italic customer quote ──────────────────────────
\newmdenv[
  topline=false,
  bottomline=false,
  rightline=false,
  linewidth=3pt,
  linecolor=QuoteBorder,
  backgroundcolor=BoxBg,
  innerleftmargin=12pt,
  innerrightmargin=12pt,
  innertopmargin=8pt,
  innerbottommargin=8pt,
  skipabove=\baselineskip,
  skipbelow=\baselineskip,
]{customerquote}

% ── spokespersonquote — indented spokesperson quote ──────────────────────────
\newmdenv[
  topline=false,
  bottomline=false,
  rightline=false,
  linewidth=3pt,
  linecolor=SectionBlue,
  backgroundcolor=BoxBg,
  innerleftmargin=12pt,
  innerrightmargin=12pt,
  innertopmargin=8pt,
  innerbottommargin=8pt,
  skipabove=\baselineskip,
  skipbelow=\baselineskip,
]{spokespersonquote}

% ══════════════════════════════════════════════════════════════════════════════
% Document
% ══════════════════════════════════════════════════════════════════════════════

\begin{document}

% ── Headline block ──────────────────────────────────────────────────────────
\begin{center}
  {\LARGE\bfseries\color{SectionBlue}
    punt-labs Launches biff: Team Communication \\
    for Engineers Who (Almost) Never Leave the Terminal} \\[0.8em]
  {\large\color{AccentGray}
    Open-source tool resurrects BSD Unix commands as MCP-native \\
    slash commands inside Claude Code --- ten commands, one install, zero context switches}
\end{center}

\vspace{0.5em}

% ── Dateline + Lede ─────────────────────────────────────────────────────────
Seattle, WA --- March 2, 2026 --- Today, punt-labs announced
\textbf{biff}, an open-source communication tool that lets software
engineers message teammates, check presence, see what their AI agents are
doing, and broadcast to their team --- all without leaving the terminal.
Unlike Slack and Discord, which demand a separate window and constant
attention, biff runs inside the engineer's existing Claude Code session as
MCP-native slash commands. Engineers type \texttt{/write @kai} the same way
they type \texttt{git push} --- in flow, with intent, and then move on.

\prsection{Problem}

Software engineers using AI-assisted coding tools like Claude Code are
experiencing a step change in productivity. Entire features that once took
days now take hours. But every time these engineers need to coordinate with
a teammate --- asking a question, sharing a diff, requesting a code review
--- they context-switch to Slack or Discord, tools designed for managers
and all-day-online knowledge workers, not for makers in deep focus. Those
tools show \emph{everything} --- every channel, every project, every server
--- regardless of which codebase is open in the terminal right now.

The cost is measurable. Gloria Mark's research at UC Irvine found that a
single interruption costs an average of 23 minutes and 15 seconds before a
worker returns to the original task. For an engineer shipping three
features a day with AI assistance, three Slack interruptions can erase an
entire feature's worth of productive time. Slack's always-on presence model
compounds the problem: even \emph{checking} for messages breaks flow, and
the social expectation of quick replies creates anxiety that degrades
concentration even when no message arrives.

The engineers building the future of software --- spending their entire day
inside a terminal, collaborating with AI agents, iterating at speeds that
would have seemed absurd two years ago --- have no communication tool that
matches how they actually work.

\prsection{Solution}

Biff brings communication into the terminal as slash commands borrowed from
the BSD Unix utilities that shipped with every workstation in the 1980s ---
\texttt{write}, \texttt{wall}, \texttt{finger}, \texttt{who}, and
\texttt{mesg} --- updated for the MCP era. Ten commands cover messaging,
presence, real-time conversation, broadcast, and availability control:

\begin{itemize}[leftmargin=1.5em, itemsep=0.2em]
  \item \texttt{/who} --- see who's online, what they're working on, and
    where each session is running
  \item \texttt{/write @kai ``ready for review''} --- send a directed
    message
  \item \texttt{/read} --- check your inbox when you're ready
  \item \texttt{/finger @kai} --- see what someone is working on without
    interrupting them
  \item \texttt{/talk @kai} --- start a real-time conversation with instant
    status bar display
  \item \texttt{/wall ``release freeze'' 2h} --- broadcast a time-limited
    announcement to the team
  \item \texttt{/plan ``refactoring auth''} --- set your status so
    teammates can see your focus
  \item \texttt{/tty main} --- name your session for multi-session work
  \item \texttt{/last} --- view session login/logout history
  \item \texttt{/mesg n} --- go do-not-disturb; \texttt{/mesg y} to
    come back
\end{itemize}

Every command implies intent. There are no channels to monitor, no threads
to catch up on, no emoji reactions to parse. The repository \emph{is} the
channel. When you open a project, you see only the teammates and messages
for \emph{that} project --- not the ten others on the back burner. Switch
repos, and your communication context switches with you. This is why biff
works for focused engineering: the same scoping that keeps your code
organized keeps your communication organized. When someone sends you a
message, your status bar shows a live indicator ---
\texttt{kai:tty1(3)} in bold yellow when you have unreads,
\texttt{kai:tty1(0)} when caught up. You decide when to engage.

Because biff speaks MCP, agents participate alongside humans as natural
members of the team. An autonomous coding agent can \texttt{/plan} what it
is working on, \texttt{/write} a human when it needs a decision, and show
up in \texttt{/who} alongside everyone else. Workflow hooks tie into the
development process: switching branches auto-sets your plan, creating a PR
triggers a \texttt{/wall} announcement. Biff is designed for hybrid teams
--- humans and agents working together --- with the same intentional,
human-scale communication vocabulary that made BSD utilities great.

% TODO: Replace with a real customer testimonial
\prsection{Customer Quote}

\begin{customerquote}
\textit{``I used to mass-quit my Slack channels once a month out of
frustration and then guiltily rejoin them. With biff, I just set
\texttt{/mesg n} when I'm deep in a problem and \texttt{/mesg y} when I
come up for air. My teammates can \texttt{/finger @priya} to see what I'm
on without interrupting. Now I run three agents alongside my own session
--- \texttt{/who} shows all four of us, each with its own plan, and when
one of them needs a decision it just \texttt{/write}s me. Last week I
shipped more in three days than I used to ship in two weeks, and I didn't
miss a single important message.''}

\raggedleft --- \textbf{Priya Chandrasekaran}, Senior Engineer at a
Series~B startup
\end{customerquote}

\prsection{Getting Started}

Getting started with biff takes two steps --- install it, then enable it
in your repo:

\begin{mdframed}[
  topline=false, bottomline=false, rightline=false, leftline=false,
  backgroundcolor=BoxBg,
  innerleftmargin=12pt, innerrightmargin=12pt,
  innertopmargin=8pt, innerbottommargin=8pt,
  skipabove=\baselineskip, skipbelow=\baselineskip,
]
\ttfamily\footnotesize
\textcolor{AccentGray}{\$} curl -fsSL https://raw.githubusercontent.com/punt-labs/biff/419ac99/install.sh | sh \\[0.6em]
\textcolor{AccentGray}{\$} biff enable
\end{mdframed}

\noindent The first command installs the CLI and registers the Claude Code
plugin. The second activates biff in the current repository and deploys
git hooks for workflow integration. Restart Claude Code, and you are live.

Biff ships with a shared relay on Synadia Cloud, so there is zero
infrastructure to provision. Team configuration lives in a \texttt{.biff}
file committed to the git repository --- the repo is both the codebase and
the communication scope. Display names are resolved automatically from
GitHub identity. When a teammate clones a repo that has a \texttt{.biff}
file, biff connects them to the right team automatically. There is no
account to create, no workspace to configure, and no browser to open. Type
\texttt{/who} to see your teammates.

\prsection{Spokesperson Quote}

\begin{spokespersonquote}
``We built biff because we watched the best engineers in the world
disappear into their terminals --- shipping at speeds that make the old
sprint-and-standup cadence look quaint --- and then lose hours every day to
communication tools designed for a different era. Slack was built for the
open-office, always-online workplace. Biff is built for the deep-focus,
AI-accelerated one. We didn't add a chat feature to the terminal. We
resurrected the Unix communication vocabulary --- \texttt{write},
\texttt{wall}, \texttt{finger}, \texttt{who}, \texttt{mesg} --- because
those commands understood something Slack forgot: communication should be
purposeful, not ambient.''

\raggedleft --- \textbf{Jim Freeman}, Founder, punt-labs
\end{spokespersonquote}

\prsection{Call to Action}

Biff is open source and available today at
\href{https://github.com/punt-labs/biff}{github.com/punt-labs/biff}.
Install with one command, restart Claude Code, and type \texttt{/who}. The
core tool and the shared relay are free under the MIT license.

% ── Boilerplate ─────────────────────────────────────────────────────────────
\vspace{2em}
\noindent\rule{\textwidth}{0.4pt}
\vspace{0.5em}

\noindent\textbf{About punt-labs}

\noindent punt-labs builds open-source developer tools for engineers who
live in the terminal. Its projects include biff (team communication),
quarry (local semantic search), and a suite of language-learning tools ---
all built as CLI applications and MCP servers that integrate natively with
AI coding assistants. More at
\href{https://github.com/punt-labs}{github.com/punt-labs}.

\end{document}
